\newcommand{\objone}{
The objective of this assignment is to develop a Python program that reads marks obtained in three different subjects and computes their average. Through this task, students will apply fundamental programming concepts such as input and output operations, variable declaration, data type handling, and arithmetic calculations. The assignment aims to strengthen understanding of basic syntax and program flow while enhancing problem-solving skills and logical reasoning in Python.
}

\newcommand{\objtwo}{
The objective of this assignment is to create a Python script that converts weight from kilograms to pounds using the standard conversion factor. This task helps in understanding user input handling, variable usage, arithmetic operations, and type conversion in Python. By implementing the formula (1 kilogram = 2.20462 pounds), students will practice writing simple computational programs and strengthen their grasp of basic programming logic and mathematical expressions.
}

\newcommand{\objthree}{
The objective of this assignment is to develop a Python program that accepts the lengths of the three sides of a prism as integer inputs and calculates its surface area using the formula ( 2ab + 2bc + 2ca ). 
}

\newcommand{\objfour}{
The objective of this assignment is to design a Python program that calculates the speed of a plane using a given distance of 395,000 meters and a travel time of 9,000 seconds. This task helps students understand the practical applicationof the speed formula ( $Speed = \frac{Distance}{Time}$). 
}

\newcommand{\objfive}{
The objective of this assignment is to develop a Python program that calculates the time required to empty a swimming pool based on its dimensions and the pump’s flow rate. The program should first compute the volume of the pool using the formula ( $\text{Volume} = \text{length} \times \text{width} \times \text{height}$ ), and then determine the time required using ( $\text{Time} = \frac{\text{Volume}}{\text{Flow Rate}}$ ). This task strengthens understanding of arithmetic operations, formula implementation, variable handling, and real-world problem modeling in Python. It also enhances logical reasoning by connecting mathematical concepts with program design.
}

\newcommand{\objsix}{
The objective of this assignment is to develop a Python script that converts temperature from Celsius to Fahrenheit, where the Celsius value is provided as a float input. The program should implement the standard conversion formula ( $C = \frac{5}{9}(F - 32)$ ). 
}

\newcommand{\objseven}{
The objective of this assignment is to develop a Python program that computes the total number of seconds in one full day. The program should apply the relationship between time units by converting hours to minutes and minutes to seconds using appropriate arithmetic operations. 
}

\newcommand{\objeight}{
The objective of this assignment is to develop a Python program that calculates the acceleration of a car using the formula ( $\frac{v_{final} - v_{initial}}{\text{time}}$ ), where the car starts from rest and reaches a velocity of 10 m/s in 20 seconds. This task helps students understand the practical application of motion equations in programming. It reinforces concepts such as variable assignment, arithmetic operations, formula implementation, and basic physics calculations in Python. Additionally, it strengthens the ability to translate mathematical expressions into structured and executable program logic.
}

\newcommand{\objnine}{
The objective of this assignment is to develop a Python program that generates and displays Pascal’s Triangle up to a specified number of rows. This task helps students understand nested loops, conditional logic, and iterative computation. It reinforces the concept of binomial coefficients, where each element in the triangle is derived from the sum of the two elements directly above it. Additionally, the assignment enhances logical thinking, pattern recognition, and structured program design in Python by translating a mathematical pattern into an algorithmic solution.
}

\newcommand{\objten}{
The objective of this assignment is to develop a Python program that displays the given numeric pattern using nested loops. This task is designed to strengthen understanding of loop control structures, particularly the use of outer and inner loops for row-wise and column-wise pattern generation. It enhances logical reasoning by requiring students to manage sequential number printing and reverse number sequences within the same row. 
}

\newcommand{\objeleven}{
The objective of this assignment is to develop a Python program that uses a while loop to calculate the sum of all even numbers between 100 and 200. This task reinforces understanding of iterative control structures, conditional statements, and arithmetic operations in Python. It helps students practice loop initialization, condition checking, and increment logic while ensuring only even numbers are included in the summation. Additionally, the assignment strengthens logical reasoning and algorithm development skills by requiring structured control flow to solve a numerical computation problem.
}

\newcommand{\objtwelve}{
The objective of this assignment is to develop a Python program that computes the sum of two mathematical series up to a given value of ( n ).

For part (a), the program should calculate the harmonic series:

$1 + \frac{1}{2} + \frac{1}{3} + \dots + \frac{1}{n}$


For part (b), the program should compute the series:

$\frac{1}{1} + \frac{2^2}{2} + \frac{3^3}{3} + \dots + \frac{n^n}{n}$
}

\newcommand{\objthirteen}{
The objective of this assignment is to develop a Python program that calculates the depreciation value of an asset by reading the purchase amount (amt), the number of years of service (year), and the rate or value of depreciation. The program should apply an appropriate depreciation formula (such as straight-line depreciation, if specified) to determine the asset’s reduced value over time. This task reinforces understanding of user input handling, arithmetic computations, formula implementation, and financial calculations in Python. It also enhances analytical thinking by applying mathematical models to real-world asset management scenarios.
}
\newcommand{\objfourteen}{
The objective of this assignment is to develop a Python program that accepts a string input from the user and displays the same string after removing all vowels. This task strengthens understanding of string manipulation, character iteration, and conditional statements in Python. It requires identifying vowel characters (both uppercase and lowercase) and filtering them out using appropriate logic. Additionally, the assignment enhances proficiency in working with loops and string processing techniques, which are fundamental concepts in programming and text-based data handling.
}
\newcommand{\objfifteen}{
The objective of this assignment is to develop a Python function that inserts a given string into the middle of another string. This task strengthens understanding of string manipulation, indexing, slicing, and function definition in Python. It requires determining the midpoint of the original string and using appropriate slicing techniques to insert the new substring accurately. Additionally, the assignment enhances problem-solving skills by applying string operations to modify and reconstruct text data in a structured and logical manner.
}
\newcommand{\objsixteen}{
The objective of this assignment is to develop a Python program that sorts the characters of a given string in lexicographical (alphabetical) order. This task reinforces understanding of string manipulation, comparison operations, and sorting techniques in Python. It requires applying appropriate methods such as built-in sorting functions or algorithmic approaches to arrange characters based on their ASCII or Unicode values.
}
\newcommand{\objseventeen}{
The objective of this assignment is to develop a Python program that replaces a specified substring within a given string with another substring without using built-in string replacement methods. This task strengthens understanding of string traversal, indexing, slicing, and manual pattern matching techniques.
}
\newcommand{\objeighteen}{
The objective of this assignment is to develop a Python program that concatenates two strings into a third string without using the `+` operator. This task reinforces understanding of string handling, iteration, and alternative concatenation techniques such as loops or string methods that do not rely on direct addition. It encourages students to manually construct the resulting string by appending characters sequentially, thereby strengthening their grasp of string traversal and memory handling. Additionally, the assignment enhances logical reasoning and deepens understanding of how string concatenation works internally in programming.
}
\newcommand{\objnineteen}{
The objective of this assignment is to develop a Python program that removes (strips) a specified set of characters from a given string. This task strengthens understanding of string manipulation, character comparison, and iterative processing in Python. It requires designing logic to identify and eliminate unwanted characters without relying solely on high-level built-in methods, thereby reinforcing fundamental programming concepts. Additionally, the assignment enhances analytical thinking by requiring the implementation of custom string filtering techniques to modify textual data effectively.
}
\newcommand{\objtwenty}{
The objective of this assignment is to develop a Python program that extracts the first n characters from a given string. This task reinforces understanding of string indexing and slicing techniques in Python. It requires handling user input, managing integer values for n, and ensuring proper boundary conditions when the specified number exceeds the string length.
}

\newcommand{\objtwentyone}{
The objective of this assignment is to develop a Python program that generates a list of 10 random integers and then separates the elements into two distinct lists: odd list and even list, based on their parity. This task reinforces understanding of list data structures, random number generation, iteration using loops, and conditional statements. It also enhances proficiency in list manipulation techniques such as element filtering and dynamic list construction. Additionally, the assignment strengthens algorithmic thinking by applying classification logic to organize data efficiently.
}

\newcommand{\objtwentytwo}{
The objective of this assignment is to develop a Python program that sorts a list of elements in ascending order using the Insertion Sort algorithm. This task strengthens understanding of fundamental sorting techniques and algorithmic logic. It requires implementing comparison and shifting operations to insert each element into its correct position within the sorted portion of the list. The assignment enhances knowledge of time complexity, iterative control structures, and in-place sorting mechanisms. Additionally, it promotes deeper insight into how basic sorting algorithms operate internally without relying on built-in sorting functions.
}

\newcommand{\objtwentythree}{
The objective of this assignment is to develop a Python program that implements the **Binary Search** algorithm to locate a specified key element within a list. This task reinforces understanding of divide-and-conquer strategy, conditional branching, and iterative or recursive logic. It requires working with a sorted list and systematically reducing the search space by comparing the key with the middle element. The assignment enhances knowledge of algorithm efficiency, particularly time complexity ( $O(\log n)$ ), and strengthens problem-solving skills by implementing an optimized searching technique without relying on built-in search functions.
}

\newcommand{\objtwentyfour}{
The objective of this assignment is to develop a Python program that creates a list containing the first eight letters of the alphabet and applies slicing operations to extract specific subsets of elements. This task reinforces understanding of list data structures and indexing mechanisms in Python. It requires using slice notation to retrieve the first three elements, select a sequence of three elements from the middle of the list, and extract elements from a specified index to the end. The assignment enhances proficiency in list manipulation, index-based access, and boundary handling while strengthening logical reasoning in data extraction using slicing techniques.
}

\newcommand{\objtwentyfive}{
The objective of this assignment is to develop a Python program that sorts a list of elements in ascending order using the **Selection Sort** algorithm. This task strengthens understanding of comparison-based sorting techniques and iterative control structures. It requires identifying the minimum element from the unsorted portion of the list and swapping it with the appropriate position in each pass. The assignment enhances knowledge of algorithm design, in-place sorting mechanisms, and time complexity analysis ( $O(n^2)$ ). Additionally, it promotes deeper comprehension of how fundamental sorting algorithms function without relying on built-in sorting methods.
}

\newcommand{\objtwentysix}{
The objective of this assignment is to develop a Python program that identifies and prints the maximum value from the second half of a given list. This task reinforces understanding of list indexing, slicing operations, and comparison techniques in Python. It requires determining the midpoint of the list, extracting the second half using slicing, and then applying logic to find the maximum element within that subset.
}

\newcommand{\objtwentyseven}{
The objective of this assignment is to develop a Python program that generates a list of numbers from 1 to 100 which are divisible by 5 or 6. This task reinforces understanding of list construction, iteration using loops, and conditional statements.
}

\newcommand{\objtwentyeight}{
The objective of this assignment is to develop a Python program that prompts the user to input five numbers and then invokes a separate function to compute their Greatest Common Divisor (GCD). This task reinforces understanding of function definition, parameter passing, and modular programming. It requires implementing an appropriate algorithm—such as the Euclidean algorithm—to calculate the GCD efficiently. The assignment enhances problem-solving skills, promotes code reusability through functional decomposition, and strengthens knowledge of number theory concepts within a programming context.
}

\newcommand{\objtwentynine}{
The objective of this assignment is to develop a Python function named `addfruit` that accepts a set of fruit names along with their corresponding prices and returns a dictionary containing the provided information. The function must implement proper validation logic to check for duplicate fruit entries and raise a `ValueError` exception if a fruit is already present in the collection. This task reinforces understanding of dictionary data structures, key-value mapping, function definition, parameter handling, and exception handling mechanisms in Python. 
}

\newcommand{\objthirty}{
The objective of this assignment is to develop a Python function that records Air Quality Index (AQI) values using the date as the key and the AQI reading as the corresponding value in a dictionary. The function should accept user input for the date and AQI, then construct and return a dictionary containing the entered information. This task reinforces understanding of dictionary data structures, key–value mapping, function design, and user input handling in Python. Additionally, it enhances skills in data organization and structured storage of time-series information using appropriate data representations.
}

\newcommand{\objthirtyone}{
The objective of this assignment is to develop a Python program that creates a dictionary containing usernames as keys and their corresponding passwords as values, using five sample entries. The task also requires demonstrating the use of the `clear()` method to remove all key–value pairs from the dictionary and the `fromkeys()` method to generate a new dictionary with specified keys and a common default value. This assignment reinforces understanding of dictionary creation, manipulation, and built-in methods in Python. It enhances proficiency in managing key–value data structures and strengthens conceptual clarity regarding dictionary operations and memory handling.
}

\newcommand{\objthirtytwo}{
The objective of this assignment is to develop Pythonic code that creates a dictionary using country names as keys and their respective capital cities as values, and then displays the entries in sorted order. This task reinforces understanding of dictionary data structures, key–value mapping, and efficient data organization. It requires applying sorting techniques—such as iterating over sorted dictionary keys—to present the information in an ordered manner. The assignment enhances proficiency in writing clean, idiomatic Python code while strengthening concepts related to data handling, sorting mechanisms, and structured output formatting.
}

\newcommand{\objthirtythree}{
The objective of this program is to create a dictionary with friends' names as keys and their phone numbers as values. The program displays the dictionary entries in sorted order and allows the user to search for a specific name. If the entered name is not found, the user can add a new name and phone number to the dictionary.
}

\newcommand{\objthirtyfour}{
The objective of this assignment is to develop a Python program that creates a dictionary using author names as keys and their corresponding ISBN numbers as values, with five sample entries. The program should also demonstrate the use of the `clear()` method to remove all key–value pairs from the dictionary and the `fromkeys()` method to create a new dictionary with specified keys and a common default value.
}
\newcommand{\objthirtyfive}{
The objective of this assignment is to develop a Python program that creates a list containing three elements and then converts that list into a tuple. This task reinforces understanding of fundamental data structures in Python, specifically the differences between mutable (list) and immutable (tuple) types. 
}
\newcommand{\objthirtysix}{
The objective of this assignment is to develop a Python program that unpacks the elements of a tuple into multiple variables. This task reinforces understanding of tuple data structures and the concept of sequence unpacking in Python. It requires assigning individual tuple elements to corresponding variables in a single statement, ensuring that the number of variables matches the number of elements. The assignment enhances proficiency in handling immutable sequences, improves code readability through structured assignment, and strengthens foundational knowledge of Python’s data manipulation techniques.
}
\newcommand{\objthirtyseven}{
The objective of this assignment is to develop a Python program that checks whether a specified item exists within a tuple. This task reinforces understanding of tuple data structures and membership testing operations in Python. It requires applying conditional statements along with the `in` operator (or equivalent logic) to verify the presence of an element. The assignment enhances knowledge of immutable sequences, improves proficiency in element lookup operations, and strengthens logical reasoning in validating data within structured collections.
}
\newcommand{\objthirtyeight}{
The objective of this assignment is to develop a Python program that unzips a list of tuples into separate individual lists. This task reinforces understanding of composite data structures, particularly lists and tuples, and their interaction in Python. It requires applying sequence unpacking techniques—such as using loops or the `zip()` function—to extract corresponding elements into distinct lists. The assignment enhances proficiency in data transformation, iteration control, and structured data handling, while strengthening conceptual clarity in managing grouped and parallel data collections programmatically.
}
\newcommand{\objthirtynine}{
The objective of this assignment is to develop a Python program that performs fundamental set operations, including intersection, union, set difference, and symmetric difference. This task reinforces understanding of the set data structure and its mathematical foundations. It requires applying built-in set operators or methods to compare and manipulate elements between two sets. The assignment enhances proficiency in handling unordered collections of unique elements, strengthens logical reasoning in relational data analysis, and deepens conceptual clarity regarding set theory operations and their implementation in Python
}
\newcommand{\objforty}{
The objective of this assignment is to develop a Python program that demonstrates the use of the `issubset()` and `issuperset()` methods in set operations. This task reinforces understanding of set relationships and hierarchical comparisons between collections of unique elements. It requires checking whether one set is completely contained within another (subset) and whether a set contains all elements of another set (superset). The assignment enhances conceptual clarity in set theory, improves proficiency in relational data validation, and strengthens logical reasoning in comparing structured datasets programmatically.
}

\newcommand{\objfortyone}{
The objective of this assignment is to develop a Python program that accepts a specified range and generates a list of tuples, where each tuple contains a number and its corresponding square. This task reinforces understanding of iteration using loops, tuple creation, and list construction. It requires applying arithmetic operations within a structured loop to compute the square of each number and store the results as ordered pairs. The assignment enhances proficiency in handling composite data structures and strengthens logical reasoning in transforming numerical ranges into structured datasets.
}
\newcommand{\objfortytwo}{
The objective of this assignment is to develop a Python program that demonstrates how to clear all elements from a set. This task reinforces understanding of the set data structure and the use of built-in methods, specifically the `clear()` method. 
}
\newcommand{\objfortythree}{
The objective of this assignment is to develop a Python program that determines the number of elements present in a set. This task reinforces understanding of the set data structure and the concept of cardinality in collections. It requires applying the appropriate built-in function to compute the length of the set.
}
\newcommand{\objfortyfour}{
The objective of this assignment is to develop a Python program that stores the latitude and longitude coordinates of a house in a tuple and displays them. This task reinforces understanding of tuple data structures as immutable sequences used to store related data. It requires representing geographical coordinates as ordered pairs and accessing tuple elements appropriately for output. The assignment enhances conceptual clarity in handling fixed data collections and strengthens proficiency in structured data representation in Python.
}
\newcommand{\objfortyfive}{
The objective of this assignment is to develop a Python program that prompts the user to enter the name of a text file, reads all the words from the file, and displays the non-duplicate words in ascending order. This task reinforces understanding of file handling operations, including opening, reading, and processing text files. It requires applying string manipulation techniques to extract words and using a set data structure to eliminate duplicates. The program must then sort the unique words before displaying them.
This assignment enhances proficiency in file I/O operations, data cleaning, sorting mechanisms, and the use of appropriate data structures for efficient data processing, while strengthening logical reasoning in handling real-world textual data.
}
\newcommand{\objfortysix}{
The objective of this assignment is to develop a Python program that determines the size of a plain text file. This task reinforces understanding of file handling and interaction with the file system in Python. It requires accessing file metadata using appropriate standard library functions to retrieve the file size in bytes.
}
\newcommand{\objfortyseven}{
The objective of this assignment is to develop a Python program that prompts the user to enter the name of a text file and then calculates the number of vowels and consonants present in the file. This task reinforces understanding of file input operations, text processing, and character-level analysis in Python. It requires reading file content, iterating through characters, distinguishing alphabetic characters from non-alphabetic symbols, and applying conditional logic to classify vowels and consonants accurately.
}
\newcommand{\objfortyeight}{
The objective of this assignment is to develop a Python program that reads and displays the first n lines of a file, where the value of n is provided by the user. This task reinforces understanding of file input operations, user interaction, and controlled iteration. It requires opening a file safely, reading its contents line by line, and ensuring that only the specified number of lines are displayed, even if the file contains more data.
The assignment enhances proficiency in file handling techniques, loop control mechanisms, and boundary condition management, while strengthening logical reasoning in processing structured textual data efficiently.
}
\newcommand{\objfortynine}{
The objective of this assignment is to develop a Python program that prompts the user to enter a filename, reads the contents of the specified file, and counts the occurrences of each letter. This task reinforces understanding of file input operations, user interaction, and character-level text processing. It requires implementing logic to read file data, filter alphabetic characters, normalize case for accurate counting, and store frequency counts using an appropriate data structure such as a dictionary.
}
\newcommand{\objfifty}{
The objective of this assignment is to develop a Python program that reads and displays the last n lines of a file, where the value of n is provided by the user. This task reinforces understanding of file handling operations, user input processing, and list manipulation techniques. It requires reading file contents efficiently and implementing logic to extract only the specified number of lines from the end of the file, while properly handling cases where n exceeds the total number of lines.
The assignment enhances proficiency in file I/O operations, indexing and slicing mechanisms, and boundary condition management, while strengthening problem-solving skills in handling real-world text processing tasks programmatically.
}

\newcommand{\objfiftyone}{
The objective of this assignment is to develop a Python program that reads two separate files and combines each line from the first file with the corresponding line from the second file. This task reinforces understanding of file handling operations, parallel iteration, and string manipulation in Python. It requires implementing logic to read both files simultaneously and merge lines in a synchronized manner, while properly handling cases where the files may differ in length.
The assignment enhances proficiency in file I/O processing, loop control mechanisms, data alignment techniques, and structured text manipulation, strengthening problem-solving skills in handling multi-source data integration programmatically.
}

\newcommand{\objfiftytwo}{
The objective of this assignment is to develop a Python program that removes newline characters from the contents of a file. This task reinforces understanding of file input and output operations, string manipulation, and text processing techniques in Python. It requires reading the file content, identifying newline (`\n`) characters, and rewriting the modified content without line breaks—either by replacing them with spaces or removing them entirely.
The assignment enhances proficiency in handling text files, performing data cleaning operations, and managing file write-back processes. It also strengthens logical reasoning in processing and transforming structured textual data efficiently.
}

\newcommand{\objfiftythree}{
The objective of this assignment is to develop a Python program that reads and displays a random line from a file. This task reinforces understanding of file handling operations, list manipulation, and the use of randomization techniques in Python. It requires reading all lines from the file, storing them in a suitable data structure, and selecting one line randomly using appropriate logic.
}

\newcommand{\objfiftyfour}{
The objective of this assignment is to develop a Python program that reads data from one CSV file and writes the same content into another CSV file. This task reinforces understanding of file handling operations and structured data processing using the CSV format. It requires implementing appropriate read and write mechanisms—such as using Python’s `csv` module—to accurately transfer rows of data while preserving their structure.
The assignment enhances proficiency in handling tabular datasets, performing file-to-file data transfer, and managing input/output streams efficiently. It also strengthens conceptual clarity in working with structured text formats commonly used in data storage and analysis applications.
}

\newcommand{\objfiftyfive}{
The objective of this assignment is to develop a Python program that identifies and matches words containing the letter 'z'. This task reinforces understanding of string processing and pattern matching techniques in Python. It may involve iterating through words in a given text or applying regular expressions to detect occurrences of the specified character within words.
The assignment enhances proficiency in text searching, conditional evaluation, and character-based filtering. It also strengthens problem-solving skills in implementing pattern recognition logic for analyzing and extracting specific information from textual data.
}

\newcommand{\objfiftysix}{
The objective of this assignment is to understand and implement string manipulation techniques in Python by processing an IP address. Specifically, the goal is to remove all leading zeros from each segment (octet) of the IP address while preserving its correct structure. This helps in improving data formatting, ensuring standard representation of IP addresses, and enhancing problem-solving skills using Python.
}
\newcommand{\objfiftyseven}{
The objective of this assignment is to learn how to extract numeric patterns from a string using Python. Specifically, the goal is to identify and retrieve all numbers consisting of digits (0–9) with lengths between 1 and 3.
}
\newcommand{\objfiftyeight}{
The objective of this assignment is to understand substring generation in Python by extracting all possible substrings from a given string. This helps in strengthening concepts of string manipulation, nested loops, and indexing, which are essential for problem-solving in programming.
}
\newcommand{\objfiftynine}{
The objective of this assignment is to learn how to extract structured information (year, month, and date) from a URL using Python. This enhances skills in string processing, pattern recognition, and data extraction, which are important in data parsing and web-related applications.
}
\newcommand{\objsixty}{
The objective of this assignment is to learn file handling in Python and perform date format transformation. Specifically, the task focuses on reading data from a file, identifying dates in yyyy-mm-dd format, and converting them into dd-mm-yyyy format. This builds skills in file processing, string manipulation, and pattern matching.
}

\newcommand{\objsixtyone}{
The objective of this assignment is to learn string processing in Python by identifying and replacing specific words within a text. Here, the goal is to replace the word "Street" with its abbreviation "St.", which helps in understanding text manipulation and pattern matching.
}
\newcommand{\objsixtytwo}{
The objective of this assignment is to understand how to extract specific patterns (words of fixed length) from a string using Python. 
}
\newcommand{\objsixtythree}{
The objective of this assignment is to understand Object-Oriented Programming (OOP) in Python by creating a class to represent a quadratic equation. The task focuses on initializing data members and implementing a method to compute the roots of the quadratic equation using mathematical formulas.
}
\newcommand{\objsixtyfour}{
The objective of this assignment is to understand inheritance in Object-Oriented Programming (OOP) using Python. It demonstrates how a base class (Student) can be extended by a derived class to process and display the marks of the top five students, improving code reusability and organization.
}
\newcommand{\objsixtyfive}{
The objective of this assignment is to understand class design in Python by creating a Library class with relevant attributes and methods. It focuses on handling data input and performing calculations (fine based on late days), thereby strengthening concepts of OOP, data encapsulation, and method implementation.
}
\newcommand{\objsixtysix}{
The objective of this assignment is to understand multiple inheritance in Python by combining two base classes (Clock and Calendar) into a derived class (CalendarClock). It demonstrates how different functionalities (time and date) can be integrated into a single class, enhancing modularity and code reuse.
}
\newcommand{\objsixtyseven}{
The objective of this assignment is to understand how to represent and manipulate polynomials using classes in Python. It focuses on implementing addition of two polynomials by organizing coefficients and performing term-wise operations, thereby strengthening concepts of OOP, lists, and data abstraction.
}

\newcommand{\objsixtyeight}{
This assignment aims to understand the structure and syntax of JSON and how data is represented using key–value pairs. It helps in learning how to create simple JSON documents and work with them. It also focuses on parsing JSON data using Python’s built-in tools. Additionally, it develops the ability to convert JSON data into Python objects for further processing. Finally, it builds basic skills for handling real-world data from files and APIs.
}

\newcommand{\objseventy}{
XML (Extensible Markup Language) is a markup language designed to store and transport data in a structured, hierarchical format using custom-defined tags. It is both human-readable and machine-readable, and is widely used for data interchange between systems.
}

\newcommand{\objseventyone}{
NumPy (Numerical Python) provides multiple functions to efficiently create arrays for numerical computation. These functions are fundamental in data science workflows because they allow controlled initialization of data structures.
}

\newcommand{\objseventytwo}{
NumPy provides powerful indexing and slicing techniques to access and manipulate array data efficiently. Integer indexing allows direct access to specific elements using their position, while array (fancy) indexing enables selection of multiple elements using a list or array of indices. Boolean indexing is particularly useful for filtering data based on conditions, returning only the elements that satisfy a given criterion. Slicing allows extraction of subarrays using a start:stop:step format, making it easy to work with ranges of data.
}

\newcommand{\objseventythree}{
A Python program can use the Pandas library to create a one-dimensional array-like object called a Series. First, a list of data values is defined in Python. This list is then passed to the `pd.Series()` function, which converts it into a structured Series object with index labels. The Series stores data in a tabular form with both values and corresponding indices. Finally, the Series is displayed using the `print()` function. This approach is commonly used in data analysis for handling and organizing one-dimensional datasets efficiently.
}

\newcommand{\objseventyfour}{
To find the eigenvalues of a ($3 \times 3$) matrix, we form the characteristic equation by subtracting ($\lambda I$) from the matrix and taking its determinant.

$\det(A - \lambda I) = 0$

This determinant gives a cubic equation in ($\lambda$). Solving it provides three eigenvalues. For diagonal matrices, eigenvalues are simply the diagonal elements. For general matrices, expansion and simplification are required. These eigenvalues are important in many applications like stability analysis and data science.
}

\newcommand{\objseventyfive}{
The objective of this program is to create and display a Pandas Data Frame using data stored in a Python dictionary. The program demonstrates how dictionary keys can be used as column names and dictionary values can be organized into rows and columns within a Data Frame. It helps in understanding the basic concepts of data manipulation and tabular data representation using the Pandas library in Python.
}
\newcommand{\objseventysix}{
The objective of this program is to find the eigenvalues of a given 3×3 matrix. Eigenvalues are important in linear algebra and help analyze the properties of matrices. The program demonstrates the use of Python libraries to perform matrix computations efficiently. It also helps in understanding concepts related to matrix transformations and characteristic equations. The result provides the eigenvalues corresponding to the given matrix.
}
\newcommand{\objseventyseven}{
The objective of this program is to compute the dot product of two vectors using the NumPy library in Python. The program demonstrates how vector operations can be performed efficiently with NumPy. It helps in understanding the concept of vector multiplication and its applications in mathematics and data analysis. The result obtained represents the scalar product of the given vectors.
}
\newcommand{\objseventyeight}{
The objective of this program is to remove duplicate rows from a Pandas DataFrame based on a specific column. The program demonstrates how to identify and eliminate redundant data entries using Pandas functions. It helps in maintaining data quality and consistency by retaining only unique values in the selected column. This operation is useful in data cleaning and preprocessing tasks before analysis.
}
\newcommand{\objseventynine}{
The objective of this program is to create a 10×10 checkerboard matrix using the NumPy library in Python. The program demonstrates the use of array creation and indexing techniques to generate alternating patterns of values. It helps in understanding matrix manipulation and efficient array operations in NumPy. The resulting matrix resembles a checkerboard with alternating elements arranged systematically.
}
\newcommand{\objeighty}{
The objective of this program is to create a scatter plot with a regression line using the Altair library in Python. The program visualizes the relationship between two variables and displays the overall trend of the data. It helps in understanding data distribution, correlation, and predictive patterns. The regression line provides a visual representation of the best-fit relationship between the variables.
}

\newcommand{\objeightyone}{
The objective of this program is to compute the covariance matrix of a dataset using the NumPy library in Python. The program demonstrates how covariance can be used to measure the relationship between different variables in a dataset. It helps in understanding the degree to which variables vary together. The resulting covariance matrix is useful for statistical analysis, data exploration, and machine learning applications.
}
\newcommand{\objeightytwo}{
The objective of this program is to group a Pandas DataFrame by two columns and compute aggregate statistics. The program demonstrates how data can be organized into groups based on multiple categories and analyzed efficiently. It helps in calculating summary measures such as count, sum, mean, minimum, and maximum values for each group. This technique is widely used in data analysis and reporting to extract meaningful insights from datasets.
}
\newcommand{\objeightythree}{
The objective of this program is to create a cumulative sum column in a Pandas DataFrame. The program demonstrates how to calculate the running total of values in a column using Pandas functions. It helps in understanding cumulative calculations and data transformation techniques. The resulting cumulative sum column is useful for trend analysis, financial calculations, and tracking progressive changes in data.
}
\newcommand{\objeightyfour}{
The objective of this program is to generate 1000 random numbers from a normal distribution and visualize their distribution using a histogram in Altair. The program demonstrates the use of random data generation and graphical representation techniques in Python. It helps in understanding the characteristics of a normal distribution and the frequency distribution of data values. The histogram provides a clear visual summary of the generated dataset.
}
\newcommand{\objeightyfive}{
The objective of this program is to write a function that standardizes a NumPy array by transforming its values to have a mean of 0 and a variance of 1. The program demonstrates the use of statistical measures such as mean and standard deviation for data normalization. It helps in preparing data for analysis and machine learning algorithms. Standardization ensures that all values are on a common scale, improving the performance of many computational methods.
}
\newcommand{\objeightysix}{
The objective of this program is to reshape a one-dimensional NumPy array into a three-dimensional array of appropriate dimensions. The program demonstrates how data can be reorganized into multiple dimensions while preserving all elements. It helps in understanding array manipulation and multidimensional data structures in NumPy. Such transformations are useful in scientific computing, data analysis, and machine learning applications.
}
\newcommand{\objeightyseven}{
The objective of this program is to merge three Pandas DataFrames based on a common key. The program demonstrates how data from multiple sources can be combined into a single DataFrame using merge operations. It helps in understanding data integration and relational data handling in Pandas. The resulting DataFrame contains consolidated information, making data analysis and processing more efficient.
}
\newcommand{\objeightyeight}{
The objective of this program is to create an interactive chart using the Altair library with a dropdown selection filter. The program demonstrates how user input can be used to dynamically filter and visualize data. It helps in understanding interactive data exploration and visualization techniques. The dropdown filter allows users to select specific categories or values, making the chart more flexible and informative.
}
\newcommand{\objeightynine}{
The objective of this program is to detect and remove outliers from a dataset using the Z-score method in Pandas. The program demonstrates how statistical techniques can be used to identify data points that significantly deviate from the mean. It helps in improving data quality and ensuring more accurate analysis. The cleaned dataset obtained after removing outliers is better suited for statistical and machine learning applications.
}
\newcommand{\objninety}{
The objective of this program is to compute the pairwise correlation matrix of a Pandas DataFrame. The program demonstrates how to measure the strength and direction of relationships between multiple variables. It helps in identifying positive, negative, or weak correlations within the dataset. The resulting correlation matrix is useful for data analysis, feature selection, and understanding variable dependencies.
}

\newcommand{\objninetyone}{
The objective of this program is to create a heatmap visualization using the Altair library in Python. The program demonstrates how data values can be represented through color intensity in a two-dimensional grid. It helps in identifying patterns, trends, and relationships within the dataset. The heatmap provides an effective and intuitive way to visualize large amounts of data for analysis and decision-making.
}
\newcommand{\objninetytwo}{
The objective of this program is to apply a custom function to each element of a NumPy array using vectorization. The program demonstrates how vectorized operations can efficiently process array elements without using explicit loops. It helps in improving computational performance and simplifying code. This technique is widely used in numerical computing, data analysis, and scientific applications.
}
\newcommand{\objninetythree}{
The objective of this program is to convert a categorical column into one-hot encoded format using the Pandas library. The program demonstrates how categorical data can be transformed into numerical representations suitable for data analysis and machine learning models. It helps in handling non-numeric features effectively by creating separate binary columns for each category. This transformation improves the compatibility of categorical data with various analytical and predictive techniques.
}
\newcommand{\objninetyfour}{
The objective of this program is to perform left, right, inner, and outer joins between two Pandas DataFrames. The program demonstrates different methods of combining data based on common columns or keys. It helps in understanding how records are matched and merged under various join operations. These techniques are essential for data integration, analysis, and managing relational datasets in Python.
}
\newcommand{\objninetyfive}{
The objective of this program is to create a rolling window standard deviation column in a Pandas DataFrame. The program demonstrates how to calculate the standard deviation of data values over a moving window of fixed size. It helps in analyzing variability and fluctuations within a dataset over time. The resulting rolling standard deviation column is useful for trend analysis, time-series studies, and statistical data processing.
}
\newcommand{\objninetysix}{
The objective of this program is to plot multiple layered charts using the Altair library with shared axes. The program demonstrates how different data visualizations can be combined into a single chart for better comparison and analysis. It helps in understanding the relationships between multiple datasets while maintaining a common scale. The resulting layered chart provides a clear and effective way to present complex data visually.
}
\newcommand{\objninetyseven}{
The objective of this program is to calculate the percentile values of a NumPy array using the NumPy library in Python. The program demonstrates how percentiles can be used to determine the relative position of data within a dataset. It helps in understanding data distribution and identifying key statistical measures. The calculated percentile values are useful for data analysis, summarization, and decision-making processes.
}
\newcommand{\objninetyeight}{
The objective of this program is to create a time-series line chart with tooltips using the Altair library in Python. The program demonstrates how temporal data can be visualized effectively through a line chart while providing additional information through interactive tooltips. It helps in analyzing trends, patterns, and changes over time. The interactive tooltips enhance data exploration by displaying detailed values when users hover over specific data points.
}
\newcommand{\objninetynine}{
The objective of this program is to split a Pandas DataFrame into training and testing sets using random sampling. The program demonstrates how data can be divided into separate subsets for model training and performance evaluation. It helps in ensuring that machine learning models are tested on unseen data to assess their generalization ability. This process is an essential step in data preprocessing and predictive modeling.
}
\newcommand{\objonehundred}{
The objective of this program is to load data from a CSV file using Pandas, clean and handle missing values, compute summary statistics using NumPy, and visualize the results with Altair. The program demonstrates the complete workflow of data preprocessing, statistical analysis, and data visualization. It helps in understanding how to prepare raw data for analysis and extract meaningful insights.
}